% ===== A 卷参考答案与评分标准 =====
\anstou{A}

\datiA{名词解释（每小题 5 分，共 20 分）}

\thao{1.}\da 物质是标志客观实在的哲学范畴，它不依赖于人的意识而存在，并能为人的意识所反映。\pingfen{3 分}物质的唯一特性是客观实在性，这一特性把物质同具体的物质形态区别开来。\pingfen{2 分}

\thao{2.}\da 度是保持事物质的稳定性的数量界限，即事物的限度、幅度和范围。\pingfen{3 分}度的两端叫关节点（临界点），超出度事物就发生质变；方法论上要求把握“适度”原则。\pingfen{2 分}

\thao{3.}\da 实践是人类能动地改造客观世界的物质活动。\pingfen{2 分}它具有直接现实性、自觉能动性和社会历史性三个基本特征，\pingfen{2 分}基本形式包括物质生产实践、社会政治实践和科学文化实践。\pingfen{1 分}

\thao{4.}\da 社会存在是社会物质生活条件的总和，包括地理环境、人口因素和物质生产方式。\pingfen{3 分}其中物质生产方式是社会发展的决定力量。\pingfen{2 分}

\datiA{简答题（每小题 8 分，共 32 分）}

\thao{1.}\yaodian
\begin{enumerate}[label=(\arabic*),leftmargin=2.2em,itemsep=2pt,topsep=2pt]
\item 哲学基本问题是思维和存在（意识和物质）的关系问题。\pingfen{2 分}
\item 第一方面是何者为本原：主张物质第一性、意识第二性的是唯物主义，反之是唯心主义。\pingfen{3 分}
\item 第二方面是二者有无同一性：承认思维能正确认识存在的是可知论，否认的是不可知论。\pingfen{3 分}
\end{enumerate}

\thao{2.}\yaodian
\begin{enumerate}[label=(\arabic*),leftmargin=2.2em,itemsep=2pt,topsep=2pt]
\item 它揭示了事物发展的源泉和动力，即事物内部对立面的统一和斗争。\pingfen{2 分}
\item 它是理解唯物辩证法其他规律和范畴的钥匙，量变质变、否定之否定都以矛盾运动为基础。\pingfen{2 分}
\item 矛盾分析法是认识事物的根本方法，提供“两点论”与“重点论”相统一的方法。\pingfen{2 分}
\item 它揭示了事物普遍联系的根本内容和变化发展的内在根据。\pingfen{2 分}
\end{enumerate}

\thao{3.}\yaodian
\begin{enumerate}[label=(\arabic*),leftmargin=2.2em,itemsep=2pt,topsep=2pt]
\item 实践是认识的来源，认识产生于改造客观世界的活动。\pingfen{2 分}
\item 实践是认识发展的动力，不断提出新问题并提供新的认识工具。\pingfen{2 分}
\item 实践是检验认识真理性的唯一标准。\pingfen{2 分}
\item 实践是认识的目的，认识最终要回到实践中指导实践。\pingfen{2 分}
\end{enumerate}

\thao{4.}\yaodian
\begin{enumerate}[label=(\arabic*),leftmargin=2.2em,itemsep=2pt,topsep=2pt]
\item 社会意识与社会存在的发展具有不完全同步性，可能超前或滞后。\pingfen{2 分}
\item 社会意识与经济发展水平具有不平衡性。\pingfen{2 分}
\item 社会意识具有历史继承性，各形式之间相互影响、相互制约。\pingfen{2 分}
\item 社会意识对社会存在具有能动的反作用，这是相对独立性最突出的表现。\pingfen{2 分}
\end{enumerate}

\datiA{论述题（每小题 15 分，共 30 分）}

\thao{1.}\yaodian
\begin{enumerate}[label=(\arabic*),leftmargin=2.2em,itemsep=2pt,topsep=2pt]
\item \textbf{原理内容}：量变是质变的必要准备，没有量的积累就没有质变；\pingfen{2 分}质变是量变的必然结果；\pingfen{2 分}质变体现并巩固量变的成果，又为新的量变开辟道路；\pingfen{2 分}量变与质变相互渗透，总量变中有部分质变，质变中有量的扩张。\pingfen{2 分}
\item \textbf{方法论}：既要重视量的积累，反对急于求成；又要抓住时机促成质变，反对畏缩不前；在实际工作中坚持适度原则。\pingfen{4 分}
\item \textbf{联系实际}：学习是知识、能力长期积累的过程，短期看不出明显进步并不等于无效；只有积累到一定程度才会出现能力层次的跃升。因此应制定阶段性目标、保持持续投入，并在条件成熟时主动寻求突破。\pingfen{3 分}
\end{enumerate}
\pingfen{论述完整、层次清楚、能结合自身实际展开的可在此范围内上浮 1～2 分。}

\thao{2.}\yaodian
\begin{enumerate}[label=(\arabic*),leftmargin=2.2em,itemsep=2pt,topsep=2pt]
\item \textbf{绝对性}：任何真理都是对客观事物及其规律的正确反映，都包含不依赖于人的客观内容；人类认识按本性能够认识无限发展的物质世界。\pingfen{4 分}
\item \textbf{相对性}：任何真理都是对客观世界一定范围、一定层次的近似正确的反映，认识有待扩展和深化。\pingfen{4 分}
\item \textbf{辩证关系}：二者相互依存、相互包含，相对之中有绝对，绝对通过相对表现出来；真理是一个由相对走向绝对的无限发展过程。\pingfen{4 分}
\item \textbf{实践意义}：既要坚持真理、反对相对主义与怀疑主义，又要发展真理、反对绝对主义和思想僵化；对待既有认识应持“坚持—发展”的态度。\pingfen{3 分}
\end{enumerate}

\datiA{材料分析题（共 18 分）}

\thao{1.}
\textbf{（1）}\pingfen{9 分}
\begin{enumerate}[label=(\arabic*),leftmargin=2.2em,itemsep=2pt,topsep=2pt]
\item 社会存在决定社会意识。该县贫困的根源在于物质生产条件落后（交通闭塞、产业单一），而非群众观念落后；因此脱贫必须从改变生产条件入手，这是唯物史观的基本要求。\pingfen{4 分}
\item 社会意识对社会存在具有能动的反作用。当地在调查研究基础上形成的发展思路属于先进的社会意识，指导群众改变生产条件，推动了脱贫进程。\pingfen{3 分}
\item 社会意识具有相对独立性，观念的转变（从“等靠要”到自主发展）往往先于或滞后于物质条件的变化，并对实践产生推动或阻碍作用。\pingfen{2 分}
\end{enumerate}

\textbf{（2）}\pingfen{9 分}
\begin{enumerate}[label=(\arabic*),leftmargin=2.2em,itemsep=2pt,topsep=2pt]
\item 人民群众是社会物质财富的创造者。产业发展、修路建冷链，都是群众劳动实践的结果。\pingfen{3 分}
\item 人民群众是社会精神财富的创造者。群众的实践经验被总结提炼，形成“找对路子”的发展思路。\pingfen{3 分}
\item 人民群众是社会变革的决定力量。返乡创业、农户广泛参与，是脱贫得以实现的决定性因素；因此必须坚持一切为了群众、一切依靠群众的群众路线。\pingfen{3 分}
\end{enumerate}
